\section{Evaluation Methodology}\label{sec:method}
\subsection{Joint evaluation protocol}
Privacy, task success, business state and cost must refer to the same intervention and evaluation population. The protocol has four steps; Table~\ref{tab:joint-checklist} is its reporting form.
\begin{enumerate}
\item \textbf{Specify observations and execution authority.} Name the hidden target, the visible links, recipient rights and any supplied boundaries. State where cover ends, which state it can access and whether its responses reach the planner.
\item \textbf{Specify the contract and tolerances.} Define the requested goal, allowed collateral transitions, and the inference, quota, token and latency tolerances, before measuring.
\item \textbf{Measure matched executions and adapted attacks.} Pair baseline and policy on the same workload. Include cover-removal and schedule-aware attacks, report the strongest tested residual inference, and take task, state and cost outcomes from those protected runs.
\item \textbf{Decide.} A policy fails if any declared tolerance is exceeded: an execution outside its state or resource contract fails even when its goal passes. Passing a finite attack suite is not a secrecy proof.
\end{enumerate}

\begin{table}[t]\centering\footnotesize
\caption{Joint evaluation checklist for one intervention and population. Attach evidence or state the missing item; goal completion and whole-state correctness are separate fields.}\label{tab:joint-checklist}
\begin{tabularx}{\columnwidth}{@{}lY@{}}\toprule
$\square$ & Hidden target, population, prior or base rate \\
$\square$ & Allowed observations and supplied boundaries \\
$\square$ & Cover termination, privileges and isolation \\
$\square$ & Planner-visible cover feedback and dispatcher \\
$\square$ & Declared tolerances: inference, goal, state, tokens, bytes, latency \\
$\square$ & Adapted attacks and largest tested residual \\
$\square$ & Original-task attempts, completion and failures \\
$\square$ & Business invariants and violations \\
$\square$ & Matched bytes, latency, tokens and quota \\\bottomrule
\end{tabularx}\end{table}


\subsection{Attacks and statistics}\label{sec:stats}
Each condition has its own attack suite, fitted on its training split. Size-only attacks are logistic regression (LR) and a 300-tree random forest (RF) on destinations and ciphertext sizes, a classifier frozen on no-cover traffic, and last- and penultimate-destination rules. Timing attacks add timing features, supervised and label-free gap removal, context-aware removal and joint decoding; Supplement S2.3 defines them. The following conventions apply to every result. \emph{Sel.}\ is the attack chosen on validation data before test scoring, and primary tests use only these choices. \emph{Max} is the largest test count in a finite suite; it is descriptive and is not an upper bound on inference. \emph{Exploratory} marks an analysis designed after test results were seen. Paired tests are exact McNemar tests with Holm correction within each prespecified family; intervals use 10,000 action-stratified fixture bootstraps and condition on the fitted attacks and designed workload.

\subsection{A floor for planner cover}\label{sec:floor}
Planner cover runs $k$ complete shadow executions beside the real one. Consider an idealized attacker who separates all $k+1$ chains perfectly and recovers each chain's label, but cannot tell which chain is real. If the real chain is exchangeable with the shadows, then given the multiset of labels the real label is $c$ with probability $n_c/(k+1)$, where $n_c$ counts the chains labeled $c$. The best guess is a modal label, which succeeds with probability
\[P^*(k)=\frac{\mathbb{E}[\max_c n_c]}{k+1}.\]
With one shadow and six uniform classes, $P^*(1)=\tfrac16+\tfrac56\cdot\tfrac12=\tfrac7{12}$. Table~\ref{tab:shadow-floor} gives exact values. Drawing shadow labels without replacement from the classes other than the real one makes all labels distinct and gives $1/(k+1)$, while keeping each shadow's marginal label uniform. $P^*(k)$ is a reference, not a bound: an attacker who learns which chain is real exceeds it, and reused real fixtures with their own cache histories make exchangeability an idealization. The right comparison for a shadow policy is therefore $P^*(k)$, not the uniform $1/6$.

\begin{table}[t]\centering\footnotesize
\caption{Pair-ambiguity floor $P^*(k)$ for $k$ shadow planners and six classes, computed exactly (\texttt{experiment/shadow\_floor.py}). \emph{iid}: shadow labels drawn uniformly; \emph{distinct}: drawn without replacement from classes other than the real one. Tokens are nominal planner multiples. Only $k=1$, iid, was run.}\label{tab:shadow-floor}
\setlength{\tabcolsep}{4pt}
\begin{tabular}{@{}rrrrl@{}}\toprule
$k$ & iid & distinct & Tokens$\times$ & Live result ($k=1$) \\\midrule
0 & 1.000 & -- & 1 & \\
1 & .583 & .500 & 2 & Max 70 / 64 of 120 \\
2 & .491 & .333 & 3 & \\
3 & .456 & .250 & 4 & \\
4 & .428 & .200 & 5 & \\\midrule
\multicolumn{5}{@{}l}{Realized $k=1$ floor on the test split: 70.5/120 (21 matching labels)} \\\bottomrule
\end{tabular}\end{table}

